\subsection{Payoff boundaries and falling log-price derivatives}
\label{sec:payoff-boundary-derivatives}

I use the falling log-price Greek operator already described by
\citet{carr2001}. I derive its concrete Gaussian boundary kernels and certify
payoff-specific constants. Neither the operator, Gaussian option pricing nor
quantum pricing of an unbounded piecewise-affine payoff is presented as a
new generic result.

Let $\lambda>0$, let $Z$ have the standard normal distribution, and write
$D=d/dm$. For a reference level $S_{\rm ref}>0$, define
\[
 P_F(m)=\mathbb E F(S_{\rm ref}e^{m+\lambda Z}),\qquad
 C_r(m)=D(D-1)\cdots(D-r+1)P_F(m),\quad r\geq1.
\]
The product factors commute because each is a polynomial in $D$.
If $s=S_{\rm ref}e^m$, repeated differentiation gives
$D(D-1)\cdots(D-r+1)=s^r\partial_s^r$. For example,
$D(s^r\partial_s^r)=r s^r\partial_s^r+s^{r+1}\partial_s^{r+1}$,
which proves the recurrence by induction. Thus $C_1=P_F'$ and
$C_2=P_F''-P_F'$.

Suppose that $F$ is continuous and piecewise affine with finitely many
positive kink levels $K_j$. Its representation is
\[
 F(s)=a+bs+\sum_j\alpha_j(s-K_j)_+.
\]
Every term has at most linear growth, so its price exists absolutely under
the Gaussian model. Differentiation after Gaussian convolution is justified
on every compact $m$ interval by a Gaussian density times an exponential
and a polynomial, all of which are integrable. Equivalently, distributional
second differentiation gives $F''=\sum_j\alpha_j\delta_{K_j}$, and
\[
 C_2(m)=\mathbb E\{S^2F''(S)\}
 =\sum_j\frac{\alpha_jK_j}{\lambda}
       \phi\!\left(\frac{m-\log(K_j/S_{\rm ref})}{\lambda}\right).
\]
Indeed, the lognormal density at $K_j$ is
$\phi((\log(K_j/S_{\rm ref})-m)/\lambda)/(\lambda K_j)$,
and multiplication by $K_j^2$ supplies the displayed coefficient.
The affine components contribute zero to $C_2$.

For a vanilla call with $S_{\rm ref}=K$,
\[
 C_2(m)=\frac K\lambda\phi(m/\lambda)\geq0,
 \qquad \|C_2\|_\infty=\frac K{\lambda\sqrt{2\pi}}.
\]
The maximum occurs at $m=0$. More generally, $\alpha_j\geq0$ makes $C_2$
nonnegative. If $g_m(t)=P_F(m+\log t)$ for $t>0$, then
\[
 g_m''(t)=t^{-2}C_2(m+\log t).
\]
Consequently, for any $U\geq0$,
\[
 g_m(1+U)-g_m(1)-Ug_m'(1)
 =U^2\int_0^1(1-z)g_m''(1+zU)\,dz\geq0.
\]
On a finite support $0\leq U\leq U_*$, a certified upper norm $M_2$
gives a bounded positive correction no larger than $M_2U_*^2/2$ before
amplitude estimation. This concerns the correction on that support. The
raw payoff can remain unbounded, and the omitted Gaussian tails require
their own moment certificate. No stochastic asset law is changed by the
algebraic subtraction of known price coefficients.

\subsection{Sharp second derivative norm for the capped call}
\label{sec:capped-second-boundary-norm}

For $K,H>0$, let
\[
 P(m)=\mathbb E\min\{(Ke^{m+\lambda Z}-K)_+,H\},\quad
 c=1+H/K,\quad\ell=\log c,\quad a=\ell/\lambda.
\]
The two payoff kink weights are $+1$ at $K$ and $-1$ at $K+H$. Therefore
\[
 C_2(m)=\frac K\lambda\{\phi(z)-c\phi(z-a)\},\qquad z=m/\lambda.
\]
This boundary expression retains cancellation between the two strikes.
The generic positive threshold-mass envelope $H\overline D_2$ remains
valid but need not be sharp for this particular cap.

Put $q(z)=\phi(z)-c\phi(z-a)$. Its derivative is
$q'(z)=-z\phi(z)+c(z-a)\phi(z-a)$. On $0\leq z\leq a$, this derivative
is strictly negative, including both endpoints. Hence no stationary point
lies in that interval. Outside it, $z$ and $z-a$ have the same sign and
$q'(z)=0$ is equivalent to
\[
 f(z)=\log\!\left(\frac z{z-a}\right)-az-\ell+a^2/2=0.
\]
There is no loss of sign when taking the logarithm: the ratio of the two
signed terms in the stationary equation is positive. Its derivative obeys
\[
 f'(z)=-\frac a{z(z-a)}-a<-a<0
 \quad\hbox{on }(-\infty,0)\cup(a,\infty).
\]
The limits on the negative component are
$f(-\infty)=+\infty$ and $f(0-)=-\infty$; on the positive component they
are $f(a+)=+\infty$ and $f(+\infty)=-\infty$. Continuity and strict
monotonicity establish exactly one root $z_-$ on the first component and
one root $z_+$ on the second.

At a stationary root,
$\phi(z)=c(z-a)\phi(z-a)/z$, so
\[
 q(z)=-ca\phi(z-a)/z.
\]
It follows that $q(z_-)>0$ and $q(z_+)<0$. Moreover, the sign of $q'$
changes from positive to negative at $z_-$, and from negative to positive
at $z_+$. Both Gaussian tails tend to zero. Thus
\[
 \|C_2\|_\infty
 =\frac K\lambda\max\{q(z_-),-q(z_+)\}.
\]
This proves the global extremum classification; a finite sample of the
price derivative would not establish it.

I enclose $\ell$ and $a$ with outward Arb arithmetic. Every retained root
has rational endpoints within its stated component, with a strictly
positive enclosure of $f$ at its left endpoint and a strictly negative
enclosure at its right endpoint. The lower endpoint of the enclosure of
$a$ is a proved decreasing-derivative margin. If an intermediate value
enclosure is ambiguous, the mean-value inequality with that margin gives
a smaller enclosing bracket, which I verify again at both endpoints.
The stationary value is then enclosed on the whole root interval, with
$\ell$ and $a$ uncertainty retained. A midpoint value provides a lower
bound on the true global norm. The maximum of the two stationary upper
bounds supplies its upper bound. I can intersect this upper bound with
other independently valid global bounds without weakening the guarantee.

\subsection{Unique global peak of the capped log-price slope}
\label{sec:capped-first-boundary-norm}

The positive first derivative can be written as
\[
 C_1(m)=P'(m)=\frac K\lambda\int_0^\ell e^x
                \phi((m-x)/\lambda)\,dx
 =Ke^{m+\lambda^2/2}
   \left\{\Phi((m+\lambda^2)/\lambda)
       -\Phi((m-\ell+\lambda^2)/\lambda)\right\}.
\]
The integral expression follows by differentiating the strike-integral
price. Completing the square proves the closed form. In particular,
$C_1(m)>0$ for every finite $m$.

Normalise the integral density on $[0,\ell]$ and call its random variable
$X$. Its density is proportional to
$\exp\{-(x-(m+\lambda^2))^2/(2\lambda^2)\}$ on that interval.
Differentiation of its normalising constant gives
\[
 \frac{d}{dm}\log C_1(m)=\frac{\mathbb E_m X-m}{\lambda^2},\qquad
 \frac{d^2}{dm^2}\log C_1(m)
 =-\frac1{\lambda^2}+\frac{\operatorname{Var}_m(X)}{\lambda^4}.
\]
The second identity follows by differentiating the tilted mean; the
exponential-family parameter is $m/\lambda^2$, so its derivative is
$\operatorname{Var}_m(X)/\lambda^2$.

For completeness, let $p_m$ be this normalised truncated-normal density,
let $\bar x=\mathbb E_m X$, and let $\mu=m+\lambda^2$. Since
$p_m'(x)=-(x-\mu)p_m(x)/\lambda^2$, integration by parts yields
\[
 -\frac{\operatorname{Var}_m(X)}{\lambda^2}
 =\int_0^\ell(x-\bar x)p_m'(x)\,dx
 =(\ell-\bar x)p_m(\ell)+\bar x p_m(0)-1.
\]
Both endpoint densities are positive and $0<\bar x<\ell$. Therefore
$\operatorname{Var}_m(X)<\lambda^2$, proving strict log-concavity.
At $m=0$ the logarithmic derivative is positive; at $m=\ell$ it is
negative. It consequently has exactly one zero $m_*\in(0,\ell)$.
Because its sign remains positive before that zero and negative after
it, $C_1(m_*)$ is the unique global maximum. Its Gaussian convolution
also tends to zero in both tails.

A rational root bracket with outward signs of
$P''/P'=(C_1+C_2)/C_1$ therefore encloses the peak. A useful explicit
monotonicity margin follows from the same integration-by-parts identity.
For $m\in[0,\ell_+]$, where $\ell_+$ is a certified upper endpoint,
both endpoint densities are at least
\[
 d_{\min}=\frac1{\lambda\sqrt{2\pi}}
       \exp\!\left\{-\frac{(\ell_++\lambda^2)^2}{2\lambda^2}\right\}.
\]
The truncation probability is at most one, so normalising cannot lower
this density bound. If $\ell_->0$ is a certified lower endpoint, then
\[
 \frac{d^2}{dm^2}\log C_1(m)
 \leq-\frac{\ell_-d_{\min}}{\lambda^2}<0.
\]
I retain an outward lower enclosure of this positive margin and verify
the final bracket signs. The peak enclosure can also be intersected with
$H/(\lambda\sqrt{2\pi})$ and $K+H$, both valid independently.

\subsection{Higher boundary derivatives and shared scalar bounds}
\label{sec:higher-boundary-derivatives}

For $r\geq2$, the recurrence is $C_{r+1}=(D-r)C_r$. Define
\[
 B_2(z)=1,\qquad
 B_{r+1}(z)=\lambda^{-1}\{B_r'(z)-zB_r(z)\}-rB_r(z).
\]
A single call boundary satisfies
\[
 C_r(m)=\frac K\lambda B_r(m/\lambda)\phi(m/\lambda).
\]
The capped boundary has the difference of this expression and its
$c$-weighted translate by $a$. Thus, with
$N_r=\|B_r\phi\|_\infty$, the following bounds are valid:
\[
 M_r^{\rm vanilla}=\frac K\lambda N_r,\qquad
 M_r^{\rm cap}\leq
 \min\left\{H\overline D_r,\frac{2K+H}{\lambda}N_r\right\}.
\]
For $r=2$ I replace the latter triangle bound by the sharper two-root
certificate. For $r=3,\ldots,6$, I isolate every real root of
$B_r'-zB_r$ with the previously checked exact rational Sturm routine,
and evaluate every stationary interval with outward Arb arithmetic.
Polynomial times Gaussian tails vanish, so these intervals exhaust the
possible nonzero extrema. Rational coefficients are available because
$\lambda$ is rational in the frozen family. This is a certified finite
application of established polynomial root isolation.

The same constants are supplied to every financial competitor. Useful
scalar bounds follow from $DC_2=C_3+2C_2$ and
$D^2C_2=C_4+5C_3+4C_2$. In particular, for
$h(m)=(1+C_2(m)/M_2)/2$,
\[
 \|h'\|_\infty\leq\frac{M_3+2M_2}{2M_2},\qquad
 \|h''\|_\infty\leq\frac{M_4+5M_3+4M_2}{2M_2}.
\]
For the capped-price scalar, $\|P'\|_\infty\leq M_1$ and
$\|P''\|_\infty\leq M_1+M_2$. For its first derivative,
$\|P'''\|_\infty\leq M_1+3M_2+M_3$.
Finally,
\[
 g_m^{(r)}(t)=t^{-r}C_r(m+\log t),\qquad t\geq1,
\]
so $M_6$ is a sixth derivative bound for the stable two-node Gaussian
quadrature in the ratio variable, with no additional cap multiplier.
A vanilla first log-price derivative is globally unbounded; its scalar
tables require a finite-domain bound such as
$Ke^{m_++\lambda^2/2}$ on $m\leq m_+$. The underlying Gaussian price
still exists because every lognormal moment is finite.

For scalar interpolation I also retain the independently valid
strike-integral inequalities
$\|P''\|_\infty\leq H\|\phi'\|_\infty/\lambda^2$ and
$\|P'''\|_\infty\leq H\|\phi''\|_\infty/\lambda^3$.
The boundary representation gives
\[
 \|C_2'\|_\infty\leq
 \frac{2K+H}{\lambda^2}\|\phi'\|_\infty,\qquad
 \|C_2''\|_\infty\leq
 \frac{2K+H}{\lambda^3}\|\phi''\|_\infty.
\]
Its strike-integral representation supplies two further bounds,
\[
 \|C_2'\|_\infty\leq
 H\left(\frac{\|\phi'\|_\infty}{\lambda^2}
       +\frac{\|\phi''\|_\infty}{\lambda^3}\right),\qquad
 \|C_2''\|_\infty\leq
 H\left(\frac{\|\phi''\|_\infty}{\lambda^3}
       +\frac{\|\phi'''\|_\infty}{\lambda^4}\right).
\]
Taking the minimum of these proved envelopes and the falling-Greek sum
bounds is valid. The record identifies each bound's origin; every method
can use the same improvement.

For a vanilla first remainder, the useful ratio derivative remains bounded
on a finite outer support even though $P'$ is globally unbounded. In fact,
\[
 g_m'(t)=t^{-1}P'(m+\log t)
 \leq t^{-1}Ke^{m+\log t+\lambda^2/2}
 =Ke^{m+\lambda^2/2}.
\]
Therefore $g_m(1+U)-g_m(1)\leq
Ke^{m+\lambda^2/2}U$. This bound justifies a first-remainder normaliser
based on the finite geometric-level maximum, rather than incorrectly
postulating a global vanilla log-price slope norm.
